\chapter{Zadání a harmonogram praxe}
\label{chap:assignment_schedule}

Tato kapitola detailně specifikuje zadané úkoly, které byly předmětem odborné praxe, a popisuje časový harmonogram jejich plnění. Práce byla rozdělena do tří hlavních oblastí: vývoj nového modulu Importér, údržba a rozvoj stávajícího frontendu (včetně implementace AI) a tvorba automatizovaných testů.

\section{Specifikace zadaných úkolů}

\subsection{Projekt Importér}
Hlavním úkolem praxe byl vývoj nového webového rozhraní pro import dat do systému Evolio jako náhrady stávající desktopové aplikace (motivace a srovnání s původním řešením jsou podrobně rozebrány v kapitole \ref{chap:importer_implementation}). Jednalo se o vývoj zcela nové komponenty bez návaznosti na existující frontendový kód.

Klíčové požadavky:
\begin{itemize}
    \item \textbf{Intuitivní UX/UI:} Důraz na jednoduchost, aby uživatelé mohli sami importovat entity jako advokátní a realitní spisy, pohledávky, subjekty, platby a finanční položky.
    \item \textbf{Implementace designu:} Kódování aplikace podle dodaných grafických návrhů s možností vlastních vylepšení.
    \item \textbf{Backend a databáze:} Spolupráce na návrhu databázových struktur a vývoj C\# doplňku (Add-in) pro zpracování dat.
\end{itemize}

\subsection{Vývoj AI rozhraní a správa frontendových požadavků}
Druhou oblastí byla práce na stávajícím frontendu systému Evolio. Nejvýznamnějším úkolem v této části byla implementace uživatelského rozhraní pro AI asistenta pro klienta z právního sektoru. Konkrétně šlo o vývoj chatovacího okna s výběrem jazykového modelu, podporou analýzy nahraných dokumentů a streamovaným zobrazením odpovědi v reálném čase. Technické detaily tohoto řešení jsou rozebrány v kapitole \ref{chap:fe_issues_ai}. Vedle tohoto hlavního úkolu jsem se průběžně podílel na řešení běžných chyb a drobných úprav v systému Evolio dle požadavků klientů.

\subsection{Automatizované testování}
Třetím pilířem praxe byla tvorba automatizovaných testů pokrývajících klíčové agendy systému Evolio (Podatelna, Fakturace, Správa případů a spisů, Evidence subjektů a plateb). Úkolem bylo rozšiřovat stávající sadu o nové scénáře a udržovat její stabilitu v rámci CI/CD procesu. Metodika, volba nástrojů a architektura testovacího projektu jsou popsány v kapitole \ref{chap:testing}.

\section{Analýza požadavků modulu Importér}
\label{sec:requirements}

Pro přehlednost jsou níže uvedeny funkční a nefunkční požadavky formulované na základě konzultací s produktovým manažerem a firemními konzultanty. Tento přehled sloužil jako vstup pro návrh datového modelu i uživatelského rozhraní. Funkční požadavky popisující chování systému z pohledu uživatele shrnuje tabulka \ref{tab:functional_requirements}, nefunkční požadavky týkající se kvality a omezení řešení pak tabulka \ref{tab:nonfunctional_requirements}.

\begin{table}[!htbp]
\centering
\caption{Funkční požadavky modulu Importér}
\label{tab:functional_requirements}
\begin{tabular}{|l|p{3.3cm}|p{7.2cm}|}
\hline
\textbf{ID} & \textbf{Název} & \textbf{Popis} \\ \hline
FR-01 & Nahrání souboru & Uživatel může nahrát vstupní soubor ve formátu Excel (XLSX) nebo CSV. \\ \hline
FR-02 & Definice šablony & Uživatel může vytvořit, upravit nebo smazat šablonu importu s uloženým mapováním sloupců. \\ \hline
FR-03 & Mapování sloupců & Systém umožňuje přiřadit sloupce vstupního souboru na pole importovaného kontextu (např. Subjekt, Adresa). \\ \hline
FR-04 & Hierarchie kontextů & Uživatel může definovat vnořené kontexty (Subjekt $\rightarrow$ Adresa $\rightarrow$ Kontakt) a více instancí téhož kontextu. \\ \hline
FR-05 & Transformace hodnot & Systém podporuje definici převodních pravidel pro textové hodnoty na interní identifikátory. \\ \hline
FR-06 & Dvoufázový import & Systém provádí nejdříve validační (přípravnou) fázi a teprve po potvrzení uživatele ostrý zápis dat. \\ \hline
FR-07 & Real-time sledování & Uživatel je informován o průběhu zpracování jednotlivých řádků v reálném čase. \\ \hline
FR-08 & Auditování a logy & Každý import je uložen včetně výsledku jednotlivých řádků pro pozdější dohledání chyb. \\ \hline
\end{tabular}
\end{table}

\begin{table}[!htbp]
\centering
\caption{Nefunkční požadavky modulu Importér}
\label{tab:nonfunctional_requirements}
\begin{tabular}{|l|p{3.3cm}|p{7.2cm}|}
\hline
\textbf{ID} & \textbf{Kategorie} & \textbf{Popis} \\ \hline
NFR-01 & Výkon & Vykreslení mapovací obrazovky se stromem do 50 kontextů musí proběhnout bez zaznamenatelné latence (pod 200 ms). \\ \hline
NFR-02 & Škálovatelnost & Systém musí zpracovat soubory obsahující řádově tisíce řádků bez blokace uživatelského rozhraní. \\ \hline
NFR-03 & Bezpečnost & Přístup k modulu je omezen pomocí oprávnění ve formě rolí systému Evolio; data v tranzitu jsou chráněna protokolem HTTPS. \\ \hline
NFR-04 & Konzistence dat & Souběžné zápisy jsou ošetřeny pomocí optimistického řízení souběhu (\texttt{SysRowVersion}). \\ \hline
NFR-05 & Rozšiřitelnost & Definice nových typů importu je možná bez zásahu do zkompilovaného kódu (pouze přes databázovou konfiguraci). \\ \hline
NFR-06 & Použitelnost & Běžný uživatel (klient) musí být schopen provést import bez asistence firemního konzultanta. \\ \hline
NFR-07 & Auditovatelnost & Každá operace zápisu zaznamenává identifikátor autora a časovou značku. \\ \hline
\end{tabular}
\end{table}
\section{Časový plán a náročnost}

Odborná praxe byla zahájena v srpnu 2025 a probíhá kontinuálně až do současnosti. Práce na jednotlivých úkolech neprobíhala striktně lineárně, ale prolínala se podle aktuálních priorit vývojového týmu.

\textbf{Přibližné rozložení časové dotace:}
\begin{itemize}
    \item \textbf{60 \%} – Vývoj modulu Importér (FE + Add-in).
    \item \textbf{20 \%} – Implementace AI rozhraní a správa frontendových požadavků.
    \item \textbf{20 \%} – Tvorba automatizovaných testů a ostatní činnosti.
\end{itemize}

Vývoj Importéru začal implementací frontendu. Práce na backendové části (C\# Add-in) byla zahájena později, jakmile byla klientská část dostatečně připravena pro integraci. Následně probíhal vývoj obou částí souběžně.